\documentclass[10pt,a4paper]{article}
\usepackage[margin=18mm]{geometry}
\usepackage{amsmath,booktabs,array,tabularx}
\usepackage{cite,microtype}
\usepackage{xcolor}
\definecolor{SrcJin}{HTML}{56617B}
\definecolor{SrcWei}{HTML}{895B32}
\definecolor{SrcNS}{HTML}{317D37}
\definecolor{SrcHan16}{HTML}{7146A9}
\definecolor{SrcHan17}{HTML}{165CA7}
\definecolor{SrcCarey}{HTML}{947112}
\definecolor{SrcThesis}{HTML}{A5418A}
\definecolor{SrcWright}{HTML}{23828C}
\definecolor{Own}{HTML}{B53D36}
\definecolor{Standard}{HTML}{2A2A2A}
\newcommand{\Jin}[1]{\textcolor{SrcJin}{#1}}
\newcommand{\Wei}[1]{\textcolor{SrcWei}{#1}}
\newcommand{\NS}[1]{\textcolor{SrcNS}{#1}}
\newcommand{\Hsix}[1]{\textcolor{SrcHan16}{#1}}
\newcommand{\Hseven}[1]{\textcolor{SrcHan17}{#1}}
\newcommand{\Carey}[1]{\textcolor{SrcCarey}{#1}}
\newcommand{\Thesis}[1]{\textcolor{SrcThesis}{#1}}
\newcommand{\Wright}[1]{\textcolor{SrcWright}{#1}}
\newcommand{\Own}[1]{\textcolor{Own}{#1}}
\newcommand{\Std}[1]{\textcolor{Standard}{#1}}
\newcommand{\scite}[2]{\textcolor{#1}{\cite{#2}}}
\newcommand{\relB}{\textsf{B}}
\newcommand{\relA}{\textbf{\textsf{A}}}
\newcommand{\relD}{\textbf{\textsf{D}}}
\newcommand{\relS}{\textsf{S}}
\renewcommand{\citepunct}{], [}
\renewcommand{\citedash}{]--[}
\input{glyphtounicode}
\pdfgentounicode=1
\usepackage[hidelinks]{hyperref}
\hypersetup{pdftitle={An established traffic backbone for an elementary VSL task},
 pdfauthor={Differentiable traffic task review},
 pdfsubject={Literature comparison and formed modeling backbone; dependency hierarchy}}
\setlength{\parindent}{0pt}
\setlength{\parskip}{4pt}
\setlength{\abovedisplayskip}{4pt}
\setlength{\belowdisplayskip}{4pt}
\setlength{\abovedisplayshortskip}{3pt}
\setlength{\belowdisplayshortskip}{3pt}
\setlength{\tabcolsep}{4pt}
\renewcommand{\arraystretch}{1.10}
\newcommand{\head}[1]{\par\vspace{5pt}\noindent\textbf{#1}\par\nobreak\vspace{1pt}}
\newenvironment{points}{\begin{list}{$\bullet$}{%
 \setlength{\leftmargin}{1.3em}\setlength{\labelwidth}{0.6em}%
 \setlength{\labelsep}{0.4em}\setlength{\itemsep}{2pt}%
 \setlength{\parsep}{0pt}\setlength{\topsep}{2pt}}}{\end{list}}
\begin{document}
\begin{center}
{\large\bfseries An established traffic backbone for an elementary VSL task}\\[3pt]
Literature comparison and formed modeling backbone --- source hierarchy, 22 October 2026
\end{center}
\vspace{-3pt}

\head{1. What a suitable backbone must provide}
Can an established variable speed limit (VSL) model support the information,
control and gradient questions with few traffic-model changes? Consider an
ending lane E, its adjacent continuing lane A and farther lane B. Initial
states can share the same E and combined A+B traffic at every road position
while differing in A/B loading. We ask whether that difference changes merge
service, useful common-speed advice and its sensitivity. A suitable backbone
must represent the distinction, conserve traffic, respect receiving space
and count every initial unit's time until exit. Simplicity concerns the
traffic rules we must author and justify, not citation count.

\head{2. What the closest published models offer}
Aggregate VSL models support questions about traffic location or history;
lane models additionally resolve A/B loading. The inspected starting points
offer the following tradeoff.

\begin{center}
\textbf{Table 1. Published starting points and their fit to the question}\par
\vspace{3pt}
\begin{tabularx}{\linewidth}{@{}>{\raggedright\arraybackslash}p{0.19\linewidth}
 >{\raggedright\arraybackslash}p{0.37\linewidth}
 >{\raggedright\arraybackslash}X@{}}
\toprule
Published model & What is already supplied & Direct finding for this review\\
\midrule
\Jin{Jin and Jin}~\scite{SrcJin}{jinjin}
& Spatial and link-average traffic; VSL entry control, capacity drop and
upstream holding queues.
& Simpler if the question concerns \emph{where} traffic lies along the
road. It changes the A/B information question.\\[3pt]
\Wei{Wei, Han and Wang}~\scite{SrcWei}{wei}
& Link traffic retaining entry/speed history, interacting speed cohorts
and capacity drop.
& A strong aggregate alternative, but no explicit A/B loading. Native
advice acts on entering traffic.\\[3pt]
\NS{Nagalur Subraveti et al.}~\scite{SrcNS}{ns}
& Successive lane-cell densities, lane-change incentives and a lane-drop
case; \Hsix{Han congestion lineage}~\scite{SrcHan16}{han16}.
& Selected lane backbone. Adopt \Hseven{Han's VSL sending restriction}~\scite{SrcHan17}{han17}
and explicitly replace movement realization.\\[3pt]
\Carey{Carey et al.}~\scite{SrcCarey}{carey}
& Multilane cell transmission, mandatory movements and receiving-allocation
alternatives.
& Supports conservative service and forward fallback. More VSL/lane-end
task closure would be needed as the main backbone.\\
\bottomrule
\end{tabularx}
\end{center}

\textbf{Few-dependency construction.} Use \NS{Subraveti's lane framework} as
the backbone and follow its adoption of \Hsix{Han's congestion approach},
completing and adapting that formulation for the \Hseven{VSL task}.
Replace movement realization with one \Carey{Carey-led adaptation};
\Wright{Wright} supports its allocation rule. The thesis clarifies the backbone.
Thus \emph{one principal framework with inherited budgets and one adapted
movement component} carries the traffic model. Our \Own{scenario and composition}
remain explicit. This concentrates independent additions within the inspected
literature; global minimality is not proved. Jin--Jin is simpler if the
information question changes to spatial traffic differences.

\textbf{Source hierarchy and relationship key.} Columns show architectural
roles; colors credit rules and references, not ownership of symbols.
\par\nobreak\vspace{2pt}
\begin{tabularx}{\linewidth}{@{}>{\raggedright\arraybackslash}X>{\raggedright\arraybackslash}X>{\raggedright\arraybackslash}X@{}}
\toprule
\textbf{Principal framework}&\textbf{Inherited budget lineage}&\textbf{Movement adaptation}\\
\NS{Subraveti article} \scite{SrcNS}{ns}&
\Hsix{Han 2016} \scite{SrcHan16}{han16} (adopted in [3])&
\Carey{Carey et al.} \scite{SrcCarey}{carey} (principal)\\
\Thesis{Subraveti thesis} \scite{SrcThesis}{thesis} (clarification)&
\Hseven{Han 2017} \scite{SrcHan17}{han17} (VSL extension)&
\Wright{Wright et al.} \scite{SrcWright}{wright} (supporting)\\
\multicolumn{3}{@{}l@{}}{Comparison only: \Jin{Jin--Jin} \scite{SrcJin}{jinjin};
\Wei{Wei et al.} \scite{SrcWei}{wei}. \Own{This study}: scenario, composition and task (\relD).}\\
\midrule
\multicolumn{3}{@{}p{\linewidth}@{}}{\relB: source form retained/re-expressed; \relA: adapted/composed from the
colored source; \Own{\relD}: defined/derived here, including member choices
(not a novelty claim); \Std{\relS}: standard algebra, conservation or calculus.}\\
\bottomrule
\end{tabularx}\par

\head{3. The formed modeling backbone}
\textbf{3.1. Physical picture, task and root assumptions.}
The \NS{lane-density architecture (\relB)} forms a
\emph{lane-resolved cell-transmission-type (CTM) model} with
\Hsix{inherited Han budgets}. Physical cells retain density and finite
storage throughout the road; there is no separate link transmission model
(LTM) link. E can change into A at upstream cell boundaries; traffic in the
final E cell must merge at the exit.
Unserved mass stays in its current cell, so congestion can extend upstream.
The causal chain is:
\begin{center}
\fbox{\NS{Lane stocks}}\ $\longrightarrow$\ \fbox{\NS{Desires} / \Hsix{budgets}}\
$\longrightarrow$\ \fbox{\Carey{Adapted movements}}\ $\longrightarrow$\ \fbox{\Std{New stocks}}
\end{center}
\Own{Member choices (\relD):} $N=10$, $\Delta x=0.1$\,km and $\Delta t=1/600$\,h (6\,s).
At step $k$, density $\rho^k_{i\ell}$ and continuous stock
$n^k_{i\ell}=\Delta x\rho^k_{i\ell}$ describe section $i$, lane
$\ell\in\{E,A,B\}$. Initialize a finite cohort independently of constant
common advice $v\in[30,60]$\,km/h, with no later arrivals. Trips end when traffic
leaves the modeled road through A or B. We assume no downstream congestion
blocks the exit. \Own{Member assumptions (\relD):} common lane flow--density parameters, fixed geometry,
old traffic/space for every movement, equal positive priorities and no extra
merging-headway multiplier. The step satisfies $u_0\Delta t,w\Delta t\leq\Delta x$;
$u_0$ is nominal free speed and $w$ the backward-wave speed magnitude.

\textbf{3.2. How much traffic can move or be received?}
\NS{Subraveti adopts Han for capacity drop}~\scite{SrcNS}{ns}; we follow that
\Hsix{inherited formulation (\relB)}~\scite{SrcHan16}{han16} with
\Hseven{nominal/VSL caps (\relB)}~\scite{SrcHan17}{han17} to give $D$ and $R$.
\Own{Our completion (\relD)} uses the full ceiling without the summary's
additional density-order switch on sending. The exact assembled variant remains ours.
\Std{FD identities are standard (\relS).} The discharge ceiling $Q(x)$ equals nominal
capacity $C$ below critical density $\rho_c$, then falls linearly to
$(1-\delta)C$ at jam density $k_J$:
\begin{align}
 C&=\frac{u_0wk_J}{u_0+w},\quad \rho_c=C/u_0,\quad
 Q(x)=\min\!\left\{\Hseven{C},\Hsix{C\left[1-\delta\frac{x-\rho_c}{k_J-\rho_c}\right]}\right\},\label{eq:ceiling}\\
 \Hseven{D_{i\ell}}&=\Hseven{\min\{v\rho_{i\ell},Q(\rho_{i-1,\ell})\}},\quad
 \Hsix{R(x,y)=\min\{Q(x),w(k_J-y),w(k_J-x)+w_d(x-y)\}}.\label{eq:budgets}
\end{align}
Here $x,y$ are the current upstream/downstream densities in the receiving
lane. The three receiving bounds express discharge capacity, downstream
space and Han's additional congestion-wave constraint. Its fixed recovery
slope is \Hsix{$w_d=Q_J/(k_J-Q_J/u_0)$, $Q_J=(1-\delta)C$}.
\Own{Member parameters (\relD):} $u_0=60$, $w=20$\,km/h, $k_J=150$\,veh/km/lane and
$\delta=0.1$; initial densities lie in $[0,k_J]$.
Our \Hsix{lane-wise budget mapping (\relA)} gives $d_{i\ell}=\Delta t D_{i\ell}$ and
$r_{ij}=\Delta t R(\rho_{ij},\rho_{i+1,j})$: all straight and diagonal
arrivals to target $j$ share that \emph{one} $r_{ij}$.
At the outlets use $r_{Nj}=\Delta t R(\rho_{Nj},0)$, $j=A,B$;
upstream boundary densities are zero with no inflow. Advice caps sending;
receiving changes through the evolving densities, not directly through $v$.

\textbf{3.3. Why and where does early merging occur?}
\NS{Retained desires (\relB)}~\scite{SrcNS}{ns} compare perceived densities
$K$, with \Thesis{current/two-cell lookahead, weights $2{:}2{:}1$ (\relB)}~\scite{SrcThesis}{thesis}.
For a nonempty origin and adjacent destination,
\begin{equation}
 \color{SrcNS}
 P_{i,\ell j}=\left[\frac{I_{i,\ell j}K_{i\ell}-K_{ij}}
 {K_{i\ell}+K_{ij}}\right]_+,\qquad
 I=1+I^r+I^{kr}+I^{coop},\qquad [x]_+=\max(x,0).\label{eq:desire}
\end{equation}
Less crowded destinations attract traffic. Within the last $D_r=0.75$\,km,
\NS{E-to-A route bias (\relB)} grows as \NS{$I^r=(1-g_i/D_r)^3$}, where
$g_i=(N-i)\Delta x$; it is zero farther upstream. This adds urgency, not
a fixed merge fraction or a hard start line: density incentives can cause
earlier E-to-A movement. \Own{Member conventions (\relD)} select mode/direction
branches for keep-right/cooperation; free A can be encouraged toward B to
make room. Missing downstream E cells are omitted and weights renormalized;
A/B lookahead extends with zero outlet density. Empty origins make no
requests. E, A and B name current lanes, so early-merged traffic becomes A.

\textbf{3.4. Which attempts succeed, and what happens at the choke?}
\Carey{Movement adaptation (\relA)}~\scite{SrcCarey}{carey} replaces native
same-cell lateral/adjusted-forward coupling. \Own{Our composition (\relD)}
separates desire from service. Requests use existing forward or adjacent
cells/outlets. At one boundary, temporarily omit $i$.
\Own{Normalize lateral desires (\relD)}, leave the rest for straight
travel, and form requests $m$ from each donor's one sending amount:
\begin{equation}
 \color{Own}
 b_{\ell j}=\frac{P_{\ell j}}{\max(1,\sum_{h\ne\ell}P_{\ell h})}\ (j\ne\ell),
 \quad b_{\ell\ell}=1-\sum_{j\ne\ell}b_{\ell j},\quad
 m_{\ell j}=d_\ell b_{\ell j}.\label{eq:requests}
\end{equation}
\Carey{R3/A3 sharing and forward fallback (\relA)}~\scite{SrcCarey}{carey}
underpin the replacement. We extend its two-stream sharing to all target
requests; \Wright{Wright's priority formulation}~\scite{SrcWright}{wright}
provides supporting justification, without importing its general junction model.
Each target serves requests at a common capped level:
\begin{equation}
 \color{SrcCarey}
 z_{\ell j}=\min(m_{\ell j},\lambda_j),\qquad
 \sum_\ell z_{\ell j}=\min\!\left(r_j,\sum_\ell m_{\ell j}\right).\label{eq:sharing}
\end{equation}
Requests below the shared level are fully served; unused entitlement returns to the others.
\Carey{Residual donor traffic gets one straight fallback $y_\ell$ (\relA)}, the minimum of
its unused sending and its own lane's unused receiving amount. Thus
\Own{$f_{\ell j}=z_{\ell j}+\mathbf1_{\{j=\ell\}}y_\ell$ (\relD)} is realized movement.
An unsuccessful E-to-A attempt can therefore continue in E if space remains.
\Own{Final-E rule (\relD): $b_{EA}=1$, $y_E=0$}; E must use A's outlet or wait.
It competes with A and any B-to-A request for the same outlet budget.
Equal priorities are not a constant realized merge proportion; demand and
receiving state determine the flow. E's use leaves less service for the
other requests, whose retained stock can restrict later receiving.
There is no separate speed-mismatch or acceleration-recovery meter.
Our normalization and three-lane sequencing remain behavioral adaptations
beyond the lane-end geometry.

\Carey{Diagonal transport (\relA)} advances one section when changing lanes.
Apply old-state flows simultaneously using \Std{standard conservation (\relS)}:
\begin{equation}
 n^{k+1}_{i\ell}=n^k_{i\ell}-\sum_j f^k_{i,\ell j}
                             +\sum_h f^k_{i-1,h\ell}.\label{eq:update}
\end{equation}
Set $f_0=0$; $f_N$ counts traffic leaving the modeled road. Matched additions/subtractions
conserve mass, while donor/receiver limits protect stock and storage.

\textbf{3.5. Complete objective, information comparison and sensitivity.}
Minimize total system time $J$, counting all initial mass until exit.
Once remaining mass is below
$\Delta x\rho_c$, all requests receive service and a unit in section $i$
has remaining mean time $(N-i+1)\Delta x/v$. Choose a fixed index $K_*$
inside that regime. With $h_{i\ell}=N-i+1$ and $\mathbf1^\top n^k$ the
total unfinished mass, the exact prefix-plus-residual objective is
\begin{equation}
 J(v)=\Delta t\sum_{k<K_*}\mathbf1^\top n^k
             +\Own{\frac{\Delta x}{v}h^\top n^{K_*}}.\label{eq:objective}
\end{equation}
\Std{Population-time accounting is standard (\relS)}; \Own{the exact residual is derived here (\relD)}.
Write the complete update as $n^{k+1}=F(n^k,v)$. Between branch switches,
\Std{$s^0=0$, $s^{k+1}=F_n^{[k]}s^k+F_v^{[k]}$ (\relS)} propagates
$s^k=\partial n^k/\partial v$ through budgets, desires, allocation and stock.
Then
\begin{equation}
 J'(v)=\Delta t\sum_{k<K_*}\mathbf1^\top s^k
       +\Own{\frac{\Delta x}{v}h^\top s^{K_*}
       -\frac{\Delta x}{v^2}h^\top n^{K_*}}.\label{eq:sensitivity}
\end{equation}
\Own{This model's full derivative (\relD)} uses \Std{standard chain rule (\relS)}.
Terms measure earlier traffic changes, residual-stock changes and
direct residual travel time. Hold a valid $K_*$ fixed locally; at native
mode or active-constraint switches a two-sided derivative can fail.

\Own{Our information/control comparison (\relD)} uses this same model and objective:
one observer sees E and A+B at each section, another sees each lane.
Admit model-aware inference from allowed histories; choose advice before
the subsequent revealing discharge. A compact observer adds the final-A
density, sufficient on the existing two-state diagnostic. Separately,
compare gradients with bounded scalar search using the same information,
and retain maximum advice as a reference. These definitions isolate what
information is needed from how the advice is optimized.
The \href{vsl_joined_candidate_contract_and_qualification.pdf}{companion formulation}
specifies all bias branches and finite-domain conventions.

\par\pdfsavepos
\write16{REVIEW-MAIN-END-PAGE: \thepage; Y-SP: \the\pdflastypos}
\makeatletter
\renewcommand{\section}{\@ifstar{\head}{\head}}
\renewcommand{\@biblabel}[1]{\ifcase#1\or\Jin{[#1]}\or\Wei{[#1]}\or\NS{[#1]}\or\Hsix{[#1]}\or\Hseven{[#1]}\or\Carey{[#1]}\or\Thesis{[#1]}\or\Wright{[#1]}\else[#1]\fi}
\makeatother
\begin{thebibliography}{8}
\setlength{\itemsep}{4pt}\setlength{\parskip}{0pt}
\bibitem{jinjin} \Jin{H.-Y. Jin and W.-L. Jin, ``Control of a lane-drop bottleneck
through variable speed limits,'' arXiv:1310.2658, 2013 (archived manuscript
dated Oct. 29, 2018). [Online]. Available:
\url{https://arxiv.org/abs/1310.2658}.}
\bibitem{wei} \Wei{L. Wei, Y. Han, and M. Wang, ``A first-order link-based flow
model with variable speed limits and capacity drops for freeway networks,''
\emph{Comput.-Aided Civ. Infrastruct. Eng.}, vol.~40, no.~21,
pp.~3200--3217, 2025, doi:
\href{https://doi.org/10.1111/mice.13492}{10.1111/mice.13492}.}
\bibitem{ns} \NS{H.~H.~S. Nagalur Subraveti, V.~L. Knoop, and B. van Arem,
``First order multi-lane traffic flow model---an incentive based macroscopic
model to represent lane change dynamics,'' \emph{Transportmetrica B, Transp.
Dyn.}, vol.~7, no.~1, pp.~1758--1779, 2019 (archived article header: 2020),
doi: \href{https://doi.org/10.1080/21680566.2019.1700846}{10.1080/21680566.2019.1700846}.}
\bibitem{han16} \Hsix{Y. Han, Y. Yuan, A. Hegyi, and S.~P. Hoogendoorn,
``New extended discrete first-order model to reproduce propagation of jam
waves,'' \emph{Transp. Res. Rec.}, vol.~2560, pp.~108--118, 2016, doi:
\href{https://doi.org/10.3141/2560-12}{10.3141/2560-12}.}
\bibitem{han17} \Hseven{Y. Han, A. Hegyi, Y. Yuan, S. Hoogendoorn,
M. Papageorgiou, and C. Roncoli, ``Resolving freeway jam waves by discrete
first-order model-based predictive control of variable speed limits,''
\emph{Transp. Res. C, Emerg. Technol.}, vol.~77, pp.~405--420, 2017, doi:
\href{https://doi.org/10.1016/j.trc.2017.02.009}{10.1016/j.trc.2017.02.009}.}
\bibitem{carey} \Carey{M. Carey, C. Balijepalli, and D. Watling, ``Extending the
cell transmission model to multiple lanes and lane-changing,''
\emph{Netw. Spat. Econ.}, vol.~15, pp.~507--535, 2015, doi:
\href{https://doi.org/10.1007/s11067-013-9193-7}{10.1007/s11067-013-9193-7}.}
\bibitem{thesis} \Thesis{H.~H.~S. Nagalur Subraveti, \emph{Lane-specific Traffic
Flow Control}, Ph.D. dissertation, Delft Univ. Technol., Delft, The
Netherlands, 2021, ch.~4, doi:
\href{https://doi.org/10.4233/uuid:e405a549-b1fe-4fa1-8090-4429612359d9}{10.4233/uuid:e405a549-b1fe-4fa1-8090-4429612359d9}.}
\bibitem{wright} \Wright{M.~A. Wright, G. Gomes, R. Horowitz, and
A.~A. Kurzhanskiy, ``On node models for high-dimensional road networks,''
\emph{Transp. Res. B, Methodol.}, vol.~105, pp.~212--234, 2017, doi:
\href{https://doi.org/10.1016/j.trb.2017.09.001}{10.1016/j.trb.2017.09.001}.}
\end{thebibliography}
\end{document}
